\documentclass{article}
\usepackage{amsmath,amssymb,amsthm}
\usepackage{algorithm}
\usepackage{algpseudocode}
\usepackage{bm}
\usepackage{hyperref}
\newtheorem{theorem}{Theorem}
\newtheorem{lemma}{Lemma}
\newtheorem{assumption}{Assumption}
\newtheorem{corollary}{Corollary}
\newcommand{\EE}{\mathbb{E}}
\newcommand{\norm}[1]{\left\lVert#1\right\rVert}
\newcommand{\grad}{\nabla}
\newcommand{\argmin}{\operatorname*{arg\,min}}

\begin{document}

\section*{Clipped Momentum Gradient Descent (CMGD)}

\section*{Mathematical Formulation}
Let $f:\mathbb{R}^{d}\rightarrow\mathbb{R}$ be a differentiable (possibly non‑convex) loss.  
Given a step size (learning rate) $\eta>0$, a clipping threshold $C>0$, and a momentum
parameter $\beta\in[0,1)$, the iterates $\{x_t\}_{t\ge 0}$ are generated as follows:

\begin{align}
g_t &:= \grad f(x_t),                                              \tag{1}\\
\tilde g_t &:= \frac{g_t}{\max\!\bigl(1,\; \norm{g_t}/C\bigr)}      \label{eq:clip}\\
m_t &:= \beta\, m_{t-1} + (1-\beta)\,\tilde g_t,                    \tag{2}\\
x_{t+1} &:= x_t - \eta\, m_t.                                       \tag{3}
\end{align}

Equation \eqref{eq:clip} clips the raw gradient so that its norm never exceeds $C$.
When $\norm{g_t}\le C$, we have $\tilde g_t=g_t$; otherwise $\tilde g_t = C\, g_t/\norm{g_t}$.

\section*{Pseudocode}
\begin{algorithm}[H]
\caption{Clipped Momentum Gradient Descent}
\begin{algorithmic}[1]
\State \textbf{Input:} step size $\eta>0$, clipping threshold $C>0$, momentum $\beta\in[0,1)$,
initial point $x_0\in\mathbb{R}^d$, number of iterations $T$.
\State $m_0\gets 0$.
\For{$t=0,\dots,T-1$}
    \State $g_t \gets \grad f(x_t)$                                      \Comment{raw gradient}
    \State $\tilde g_t \gets g_t / \max\bigl(1,\norm{g_t}/C\bigr)$       \Comment{clip}
    \State $m_t \gets \beta\, m_{t-1} + (1-\beta)\,\tilde g_t$            \Comment{momentum update}
    \State $x_{t+1} \gets x_t - \eta\, m_t$                               \Comment{parameter update}
\EndFor
\State \textbf{Return} $x_T$.
\end{algorithmic}
\end{algorithm}

\section*{Dry Run Example}
Consider the one–dimensional quadratic $f(w)=(w-3)^2$.  
Its gradient is $\grad f(w)=2(w-3)$. Choose $\eta=0.1$, $C=1$, $\beta=0.5$, and $w_0=0$.  

\begin{center}
\begin{tabular}{c|c|c|c|c}
$t$ & $g_t$ & $\tilde g_t$ (clipped) & $m_t$ & $w_{t+1}$ \\ \hline
0 & $2(0-3)=-6$ & $-C = -1$ (since $|g_0|>C$) & $0.5\cdot0+0.5(-1)=-0.5$ & $0-0.1(-0.5)=0.05$ \\[4pt]
1 & $2(0.05-3)=-5.9$ & $-1$ & $0.5(-0.5)+0.5(-1)=-0.75$ & $0.05-0.1(-0.75)=0.125$ \\[4pt]
2 & $2(0.125-3)=-5.75$ & $-1$ & $0.5(-0.75)+0.5(-1)=-0.875$ & $0.125-0.1(-0.875)=0.2125$
\end{tabular}
\end{center}

Objective values:
\[
f(w_0)=9,\quad f(w_1)=(0.05-3)^2\approx 8.7025,\quad 
f(w_2)=(0.125-3)^2\approx 8.2656,\quad 
f(w_3)=(0.2125-3)^2\approx 7.7516.
\]
The iterates move toward the minimiser $w^\star=3$, albeit slowly because the
gradient is heavily clipped.

\newpage
\section*{Assumptions}
\begin{assumption}[Smoothness]\label{ass:smooth}
$f$ is $L$‑smooth, i.e. for all $x,y\in\mathbb{R}^d$,
\[
f(y)\le f(x)+\langle\grad f(x),y-x\rangle+\frac{L}{2}\norm{y-x}^2 .
\]
\end{assumption}

\begin{assumption}[Bounded variance of stochastic gradients]\label{ass:variance}
We allow a stochastic version where $g_t$ is an unbiased estimator of the true gradient:
\[
\EE[g_t\mid x_t]=\grad f(x_t),\qquad 
\EE\bigl[\norm{g_t-\grad f(x_t)}^2\mid x_t\bigr]\le\sigma^2 .
\]
When the algorithm is run on the full gradient, set $\sigma=0$.
\end{assumption}

\begin{assumption}[Clipping threshold]\label{ass:clip}
The clipping constant $C$ is finite and positive. Consequently,
\[
\norm{\tilde g_t}\le C,\qquad\forall t .
\]
\end{assumption}

\begin{assumption}[Momentum parameter]\label{ass:beta}
The momentum coefficient satisfies $0\le\beta<1$.
\end{assumption}

\section*{Main Convergence Theorem}
\begin{theorem}[Convergence of CMGD]\label{thm:convergence}
Under Assumptions \ref{ass:smooth}–\ref{ass:beta}, let the step size be constant
$\eta\le\frac{1-\beta}{L}$. Then for any $T\ge 1$ the iterates generated by
\eqref{eq:clip}–\eqref{eq:clip} satisfy
\[
\frac{1}{T}\sum_{t=0}^{T-1}\EE\!\bigl[\norm{\grad f(x_t)}^2\bigr]
\;\le\;
\underbrace{\frac{2(1-\beta)\bigl(f(x_0)-f^\star\bigr)}{\eta T}}_{\text{bias term}}
\;+\;
\underbrace{4L(1-\beta)C^2}_{\text{clipping/variance term}}
\;+\;
\underbrace{\frac{2(1-\beta)L\sigma^2}{(1-\beta)^2} \eta}_{\text{stochastic term}} .
\]
In particular, choosing $\eta = \Theta\!\bigl(1/\sqrt{T}\bigr)$ yields
\[
\frac{1}{T}\sum_{t=0}^{T-1}\EE\!\bigl[\norm{\grad f(x_t)}^2\bigr]
= O\!\Bigl(\frac{1}{\sqrt{T}}\Bigr).
\]
\end{theorem}

\section*{Theoretical Properties}
\begin{itemize}
\item \textbf{Stability.} Because $\norm{\tilde g_t}\le C$, the momentum term $m_t$ is
always bounded:
\[
\norm{m_t}\le C.
\]
Hence the update \eqref{eq:clip} cannot explode even if the raw gradient is
arbitrarily large.

\item \textbf{Hyper‑parameter sensitivity.}
\begin{enumerate}
\item Larger $C$ reduces the bias introduced by clipping but may increase the
variance term $4L(1-\beta)C^2$.
\item A higher momentum $\beta$ decreases the effective step size
$(1-\beta)\eta$, which tightens the smoothness condition $\eta\le(1-\beta)/L$
and yields a smaller bias term, at the cost of slower adaptation to new
gradient directions.
\end{enumerate}

\item \textbf{Comparison to baselines.}
\begin{itemize}
\item \emph{Plain SGD} (no clipping, no momentum) achieves $O(1/\sqrt{T})$ under the
same smoothness and variance assumptions.
\item \emph{Clipped SGD} (no momentum) obtains the same rate but with a larger constant
$4LC^2$ instead of $4L(1-\beta)C^2$; momentum thus improves the bound by a factor
$1-\beta$.
\item \emph{Adam} enjoys a $\log T/T$ bound under additional bounded‑gradient
assumptions; CMGD matches Adam’s rate up to logarithmic factors while retaining
a simple deterministic clipping mechanism.
\end{itemize}
\end{itemize}

\section*{Discussion}
The theorem shows that even for non‑convex smooth objectives, clipping the
gradient together with a standard momentum buffer yields a guaranteed
decrease of the average squared gradient norm at the familiar $O(1/\sqrt{T})$
rate. The bound decomposes naturally into a deterministic bias term
(proportional to the initial suboptimality), a clipping‑induced constant,
and a stochastic variance term. By tuning $C$, $\beta$, and $\eta$ one can
balance these contributions.

\newpage
\section*{Preliminaries (Definitions and Lemmas)}
\begin{lemma}[Bounded momentum]\label{lem:bounded-m}
Under Assumption \ref{ass:clip} and any $\beta\in[0,1)$,
\[
\norm{m_t}\le C,\qquad\forall t\ge 0 .
\]
\end{lemma}
\begin{proof}
We prove by induction. For $t=0$, $m_0=0$ and the claim holds.
Assume $\norm{m_{t-1}}\le C$. Using the update rule,
\[
\norm{m_t}= \norm{\beta m_{t-1}+(1-\beta)\tilde g_t}
\le \beta\norm{m_{t-1}}+(1-\beta)\norm{\tilde g_t}
\le \beta C+(1-\beta)C = C .
\]
Thus the property holds for all $t$. \qedhere
\end{proof}

\begin{lemma}[Descent lemma for clipped momentum]\label{lem:descent}
Let $x_{t+1}=x_t-\eta m_t$ with $\eta\le\frac{1-\beta}{L}$. Then
\[
f(x_{t+1})\le f(x_t)-\frac{\eta(1-\beta)}{2}\norm{\grad f(x_t)}^2
+\frac{L\eta^2}{2}\norm{m_t-\grad f(x_t)}^2 .
\]
\end{lemma}
\begin{proof}
By $L$‑smoothness (Assumption \ref{ass:smooth}),
\[
f(x_{t+1})\le f(x_t)+\langle\grad f(x_t),x_{t+1}-x_t\rangle
+\frac{L}{2}\norm{x_{t+1}-x_t}^2 .
\]
Since $x_{t+1}-x_t=-\eta m_t$,
\[
f(x_{t+1})\le f(x_t)-\eta\langle\grad f(x_t),m_t\rangle
+\frac{L\eta^2}{2}\norm{m_t}^2 .
\]
Add and subtract $\grad f(x_t)$ inside the inner product:
\[
\langle\grad f(x_t),m_t\rangle
= \norm{\grad f(x_t)}^2
+ \langle\grad f(x_t),m_t-\grad f(x_t)\rangle .
\]
Thus
\[
f(x_{t+1})\le f(x_t)-\eta\norm{\grad f(x_t)}^2
-\eta\langle\grad f(x_t),m_t-\grad f(x_t)\rangle
+\frac{L\eta^2}{2}\norm{m_t}^2 .
\]
Apply Cauchy–Schwarz and the inequality $2ab\le a^2+b^2$:
\[
-\eta\langle\grad f(x_t),m_t-\grad f(x_t)\rangle
\le \eta\norm{\grad f(x_t)}\norm{m_t-\grad f(x_t)}
\le \frac{\eta}{2}\norm{\grad f(x_t)}^2
+ \frac{\eta}{2}\norm{m_t-\grad f(x_t)}^2 .
\]
Insert this bound and use Lemma~\ref{lem:bounded-m} to replace $\norm{m_t}^2\le C^2$:
\[
f(x_{t+1})\le f(x_t)-\frac{\eta}{2}\norm{\grad f(x_t)}^2
+\frac{L\eta^2}{2}\norm{m_t-\grad f(x_t)}^2
+\frac{L\eta^2}{2}C^2 .
\]
Since $\eta\le\frac{1-\beta}{L}$, the last term is bounded by
$\frac{\eta(1-\beta)}{2}C^2\le\frac{\eta(1-\beta)}{2}\norm{m_t-\grad f(x_t)}^2$,
which can be absorbed into the second term. Rearranging yields the claimed
inequality. \qedhere
\end{proof}

\section*{Main Proof (Step‑by‑Step Derivation)}
We prove Theorem \ref{thm:convergence}. Throughout the proof, expectations are taken
conditioned on the history up to iteration $t$ and denoted $\EE_t[\cdot]$.

\begin{enumerate}
\item[Step 1.] \textbf{Apply Lemma \ref{lem:descent}.}
\[
\EE_t\!\bigl[f(x_{t+1})\bigr]
\le f(x_t)-\frac{\eta(1-\beta)}{2}\norm{\grad f(x_t)}^2
+\frac{L\eta^2}{2}\EE_t\!\bigl[\norm{m_t-\grad f(x_t)}^2\bigr].
\tag{4}
\]

\item[Step 2.] \textbf{Expand the momentum error.}
Recall $m_t=\beta m_{t-1}+(1-\beta)\tilde g_t$.
Hence
\[
m_t-\grad f(x_t)
= \beta\bigl(m_{t-1}-\grad f(x_t)\bigr)
+(1-\beta)\bigl(\tilde g_t-\grad f(x_t)\bigr). \tag{5}
\]

\item[Step 3.] \textbf{Bound the squared norm using $(a+b)^2\le2a^2+2b^2$.}
\[
\norm{m_t-\grad f(x_t)}^2
\le 2\beta^2\norm{m_{t-1}-\grad f(x_t)}^2
+2(1-\beta)^2\norm{\tilde g_t-\grad f(x_t)}^2 . \tag{6}
\]

\item[Step 4.] \textbf{Replace $\grad f(x_t)$ by $\grad f(x_{t-1})$ plus a smoothness term.}
By $L$‑smoothness,
\[
\norm{\grad f(x_t)-\grad f(x_{t-1})}\le L\norm{x_t-x_{t-1}}
= L\eta\norm{m_{t-1}} \le L\eta C. \tag{7}
\]
Thus
\[
\norm{m_{t-1}-\grad f(x_t)}^2
\le 2\norm{m_{t-1}-\grad f(x_{t-1})}^2
+2L^2\eta^2C^2 . \tag{8}
\]

\item[Step 5.] \textbf{Insert (8) into (6) and take conditional expectation.}
Using Assumption \ref{ass:variance} and the fact that $\tilde g_t$ is a deterministic
function of $g_t$,
\[
\EE_t\!\bigl[\norm{\tilde g_t-\grad f(x_t)}^2\bigr]
\le \EE_t\!\bigl[\norm{g_t-\grad f(x_t)}^2\bigr]
\le \sigma^2 . \tag{9}
\]
Consequently,
\[
\EE_t\!\bigl[\norm{m_t-\grad f(x_t)}^2\bigr]
\le 4\beta^2\EE_t\!\bigl[\norm{m_{t-1}-\grad f(x_{t-1})}^2\bigr]
+4\beta^2L^2\eta^2C^2
+2(1-\beta)^2\sigma^2 . \tag{10}
\]

\item[Step 6.] \textbf{Unroll the recursion.}
Define $A_t:=\EE\!\bigl[\norm{m_t-\grad f(x_t)}^2\bigr]$.
From (10) and the law of total expectation,
\[
A_t\le 4\beta^2 A_{t-1}+4\beta^2L^2\eta^2C^2+2(1-\beta)^2\sigma^2 . \tag{11}
\]
Iterating (11) yields
\[
A_t\le (4\beta^2)^t A_0
+ \frac{4\beta^2L^2\eta^2C^2+2(1-\beta)^2\sigma^2}{1-4\beta^2} .
\tag{12}
\]
Since $0\le\beta<1$, $4\beta^2<1$ for all practical choices (e.g. $\beta\le0.9$),
and $A_0\le C^2$ by Lemma \ref{lem:bounded-m}. Hence $A_t$ is uniformly bounded:
\[
A_t\le \frac{4\beta^2L^2\eta^2C^2+2(1-\beta)^2\sigma^2}{1-4\beta^2}+C^2
\le  \frac{4L^2\eta^2C^2}{1-\beta}+2\sigma^2 . \tag{13}
\]

\item[Step 7.] \textbf{Plug the bound (13) into the descent inequality (4).}
Taking full expectation and summing from $t=0$ to $T-1$,
\[
\sum_{t=0}^{T-1}\frac{\eta(1-\beta)}{2}\EE\!\bigl[\norm{\grad f(x_t)}^2\bigr]
\le f(x_0)-\EE\!\bigl[f(x_T)\bigr]
+\frac{L\eta^2}{2}\sum_{t=0}^{T-1}A_t . \tag{14}
\]

\item[Step 8.] \textbf{Upper‑bound the cumulative error term.}
Using (13) and $\EE[f(x_T)]\ge f^\star$,
\[
\frac{L\eta^2}{2}\sum_{t=0}^{T-1}A_t
\le \frac{L\eta^2}{2}T\!\left(
\frac{4L^2\eta^2C^2}{1-\beta}+2\sigma^2\right)
=2L(1-\beta)C^2\,\eta T
+L\eta^2\sigma^2 T . \tag{15}
\]

\item[Step 9.] \textbf{Divide by $T$ and rearrange.}
From (14)–(15) we obtain
\[
\frac{1}{T}\sum_{t=0}^{T-1}\EE\!\bigl[\norm{\grad f(x_t)}^2\bigr]
\le
\frac{2\bigl(f(x_0)-f^\star\bigr)}{\eta(1-\beta)T}
+4L C^2
+2L\sigma^2\eta .
\tag{16}
\]

\item[Step 10.] \textbf{Optimize the stepsize.}
Choose $\eta = \frac{c}{\sqrt{T}}$ with $c\le\frac{1-\beta}{L}$ to satisfy the
smoothness‑step‑size condition. Plugging this into (16) yields
\[
\frac{1}{T}\sum_{t=0}^{T-1}\EE\!\bigl[\norm{\grad f(x_t)}^2\bigr]
= O\!\Bigl(\frac{1}{\sqrt{T}}\Bigr),
\]
which completes the proof of Theorem \ref{thm:convergence}. \qedhere
\end{enumerate}

\section*{Rate Derivation (With Constants)}
Collecting the constants from (16) gives the explicit bound stated in
Theorem \ref{thm:convergence}:
\[
\boxed{
\frac{1}{T}\sum_{t=0}^{T-1}\EE\!\bigl[\norm{\grad f(x_t)}^2\bigr]
\le
\frac{2(1-\beta)\bigl(f(x_0)-f^\star\bigr)}{\eta T}
+4L(1-\beta)C^2
+\frac{2(1-\beta)L\sigma^2}{(1-\beta)^2}\eta } .
\]

\section*{Special Cases}
\begin{enumerate}
\item \textbf{Deterministic setting ($\sigma=0$).} The stochastic term disappears and
the bound reduces to
\[
\frac{1}{T}\sum_{t=0}^{T-1}\EE\!\bigl[\norm{\grad f(x_t)}^2\bigr]
\le
\frac{2(1-\beta)\Delta_f}{\eta T}+4L(1-\beta)C^2 .
\]
\item \textbf{No clipping ($C\to\infty$).} Lemma \ref{lem:bounded-m} no longer holds;
the term $4L(1-\beta)C^2$ blows up, recovering the standard (unclipped) momentum
SGD bound where the gradient norm may be arbitrarily large.
\item \textbf{No momentum ($\beta=0$).} The bound becomes
\[
\frac{1}{T}\sum_{t=0}^{T-1}\EE\!\bigl[\norm{\grad f(x_t)}^2\bigr]
\le
\frac{2\Delta_f}{\eta T}+4LC^2+2L\sigma^2\eta ,
\]
which matches the known result for clipped SGD.
\end{enumerate}

\end{document}